% Example LaTeX document for GP111 - note % sign indicates a comment
\documentclass[12pt]{amsart}
% Default margins are too wide all the way around. I reset them here
\setlength{\hoffset}{-.75in}
\setlength{\textwidth}{6.5in}
\setlength{\voffset}{-.5in}
\setlength{\textheight}{9.0in}

\usepackage{enumerate}
\usepackage{amsmath}
\usepackage{amssymb}
\usepackage{amsthm}
\usepackage{mathrsfs}
\usepackage{amsrefs}
\usepackage{setspace}
\usepackage{tikz}
\usepackage{graphicx}

\bibliographystyle{amsplain}

\newtheorem{thm}{Theorem}[section]
\newtheorem{cor}[thm]{Corollary}
\newtheorem{lem}[thm]{Lemma}
\newtheorem*{claim}{Claim}
\theoremstyle{definition}
\newtheorem{defn}[thm]{Definition}
\newtheorem{rmrk}[thm]{Remark}
\newtheorem{ex}[thm]{Example}
\theoremstyle{plain}

\newtheoremstyle{nameit}%
{\medskipamount}%
{\medskipamount}%
{\itshape}%
{}%
{\bfseries}%
{.}%
{\labelsep}%
{\thmnote{#3}}%

\theoremstyle{nameit}
\newtheorem{namethm}{}

\theoremstyle{plain}


\title{On The Continuum Hypothesis and the Boolean Satisfaction Problem}
\author{C. Robinson Tompkins}
\date{September 10, 2020}

\newtheorem{definition}{Definition}
%\newcommand{\deg}{\mathrm{deg}}
%\newcommand{\x}{\mathbf{x}}

\begin{document}
\maketitle
\begin{onehalfspace}

%%%%%%%%%%%%%%%%%%%%%%Introduction%%%%%%%%%%%%%%%%%%%%%%%%%%%%%%%%%%%%%%%
\section{Introduction}
In St. Petersburg in 1845, Georg Cantor was born, and with him came tremendous changes in the field
of Mathematics. Previous to Cantor's work on the infinite, it was deeply established that infinite
treatments of numbers only be concerned with treatment of arbitrarily large finite numbers. It is
even documented that Carl Friedrich Gauss (1777-1855) wrote in a letter to a friend \cite{maor}:

\begin{quotation}
I must protest something vehemently against your use of the infinte as something consummated, as
this is never permitted in mathematics. The infinte is but a \emph{fa\c{c}on de parler}...
\end{quotation}

Thus in historical writings pertaining to limits and continuity we never see $n = \infty$ and instead see
$n\to\infty$.

Cantor's work all henged on the \emph{cardinality} of sets, or more specifically their "size." Clearly, in
the Gaussian context any set with a finite number of elements is permitted, and the \emph{cardinality} of said
set is precisely the number of elements. Furthermore, it is widely accepted to write the cardinality of a set $X$
in the fashion $\vert X\vert$. An example of the \emph{cardinality} of the set of the first three natural numbers
follows as $\vert\{0, 1, 2 \}\vert = 3$.

A set $X$ is said to be of \emph{infinite cardinality} when there exists a set $Y \subseteq X$ and a mapping
$f : X \to Y$ such that $f$ is a bijection. For example again consider the natural numbers
$\mathbb{N} = \{0, 1, 2, \ldots \}$, and consider the even natural numbers defined as
$2\mathbb{N} = \{0, 2, 4, \ldots\}$. We can find such a bijection $f: \mathbb{N}\to 2\mathbb{N}$ given
as $f(n) = 2n$, for $n\in\mathbb{N}$. Thus we have that the natural numbers are indeed of \emph{infinite cardinality},
and more specifically we conventionally write:
\begin{equation}
\vert\mathbb{N}\vert=\omega.
\label{eq:cardNat}
\end{equation}

\subsection{Introducting the Continuum Hypothesis}
In (\ref{eq:cardNat}) we have the infinite cardinality of the naturals notated. For the sake of discussion we now wish
to consider the Kleene closure as definied below.
\begin{defn}[Kleene Closure]
For some alphabet $\Sigma$, $\Sigma^*$ is the Kleene closure of our alphabet. That is, $\Sigma^*$ is all
of the possible words over the alphabet $\Sigma$. Notice that $\Sigma^*$ is the free monoid over the set
$\Sigma$.
\end{defn}

Now for the sake of exploring the cardinality of the continuum $c$ or $\mathbb{R}$, consider the binary
expansions over the set $(0, 1) \subset \mathbb{R}$, which analogous to considering the Kleene closure
over the two element set, ${0,1}^*$. For some $w\in{0 ,1}^*$ we quickly see the binary expansion as $0.w$.
Furthermore, we, for convenience, occaisionally write $w=w_0 w_1 w_2 \ldots w_n \ldots w_\omega$. For finiteary
expansions we simply let $w_{n+1} = w_{n+2} = w_{n+3} = \cdots = 0$.

\begin{thm}
$\vert\mathbb{R}\vert = \vert(0, 1)\vert = \vert\{0,1\}^*\vert > |\mathbb{N}|$, or in plain language the real numbers
$\mathbb{R}$ are larger in cardinality than the natural numbers $\mathbb{N}$.
\end{thm}

\begin{proof}
We begin by assuming that we can list all of possible words $w_n \in \{0, 1 \}^*$.
\end{proof}
%%%%%%%%%%%%%%%%%%%%%%%%%%%%%%%%%%%%%%%%%%%%%%%%%%%%%%%%%%%%%%%%%%%%%%%%%%

%%%%%%%%%%%%%%%%%%%%%Preliminaries%%%%%%%%%%%%%%%%%%%%%%%%%%%%%%%%%%%%%%%%
\section{Preliminaries}

We will use the standard definitions from the Lothaire book on combinatorics on words for our definitions
with a few additions \cite{lothaire}.

We begin by defining an alphabet $\Sigma$ to be a finite set of symbols from which we will make words
by concatenation (note, we will use capital Greek letters for alphabets). Further, we define a word $W$ to
be a list of symbols from any alphabet $\Sigma$ written horizontally (we will use capital letters to denote
words and lower case letters to denote letters). We will denote the word with no letters, that is the empty
word, by $\varepsilon$.

\subsection{Words}

The length (or number of letters) for a word $W$ will be written $|W|$. Note that we will use the same
symbol to represent the size of a set or absolute value. The difference will be clear based on context.
Notice that $|\varepsilon| = 0$. Further we will represent $|W|_a$ to represent the number of times the
letter $a$ occurs in $W$. Also we will use $|W|_{aba}$ to represent the number of times the word $aba$
occurs in $W$. For example if $C = abaababa$, then we have $|C|_{aba} = 3$ along with $|C| = 8$.

A word $U$ is a factor of a word $V$ if there exist two (possibly empty) words $S$ and $T$ such that $V =
TUS$. We will also say that $U$ is a subword of $V$ (or $V$ contains $U$). If $T=\varepsilon$, then we
call $U$ the prefix of $V$. Similarly, if $S = \varepsilon$, then we call $U$ the suffix of $V$.

For some alphabet $\Sigma$, $\Sigma^*$ is the Kleene closure of our alphabet. That is, $\Sigma^*$ is all
of the possible words over the alphabet $\Sigma$. Notice that $\Sigma^*$ is the free monoid over the set
$\Sigma$.

\subsection{Morphisms}

A morphism $h$ is a mapping from $\Sigma^*$ into $\Delta^*$, where $\Sigma$ and $\Delta$ are
alphabets, such that $h(WV) = h(W)h(V)$ for all
words $W,V\in\Sigma^*$, and $h(\varepsilon) = \varepsilon$. Note that $W$ and $V$ could
potentially be single letters. We note that if $X\subseteq\Sigma$ ($X$ represents a set of words) for some
alphabet $\Sigma$, $h(X)$ represents the set of words $\{h(W):W\in X\}$. Further, we call $h$
non-erasing if for all $a\in\Sigma$, where $\Sigma$ is an alphabet, $h(a)\neq\varepsilon$.

Recall from earlier that the Thue-Morse morphism, $\mu$ defined as

\vspace{\baselineskip}
\begin{singlespace}
\begin{displaymath}
\mu(t) =
\begin{cases}
01, & \textrm{for } t=0\\
10, & \textrm{for } t=1,
\end{cases}
\end{displaymath}
\end{singlespace}
\vspace{\baselineskip}

\noindent
is a morphism defined on the alphabet with two letters. For convenience, we will call the alphabet with $n
$ letters $\Sigma_n$. Infinite words are possible with such a morphism. We have displayed the $n^
\mathrm{th}$ Thue-Morse word as being $\mu^n(0)$. We will use $\omega$ to represent the first infinite
ordinal. So the Thue-Morse infinite word becomes
\[\mathbf{T} = \lim_{n\to\infty}\mu^n(0) = \mu^\omega (0),\]
 as previously seen. Note that we will use bold capitol letters to represent infinite words, with $\mathbf{T}$
here representing
the Thue-Morse infinite word.

When discussing $\Sigma_2 = \{0,1\}$, the two letter alphabet we will use $\bar{0}$ to denote
the complement of 0 (or 1 if we need that complement). That is $\bar{0}=1$ and $\bar{1}=0$. This will
become necessary in Chapter 2.

For some morphism $h:\Sigma^*\to\Delta^*$, we will call $h$ uniform if $|h(a)| = n$ for some integer $n$
for all $a\in\Sigma$ (more exactly, in this case we will call $h$ $n$-uniform).  We will call a morphism $h:
\Sigma^*\to\Delta^*$ square-free when $h(W)$ is square-free if and only if $X\in\Sigma^*$ is square-free.
Similarly, we will call $h:\Sigma^*\to\Delta^*$ an overlap-free morphism when $h(W)$ is overlap-free if
an only if $W\in\Sigma^*$ is overlap-free.

%%%%%%%%%%%%%%%%%%%%%%%%%%%%%%%%%%%%%%%%%%%%%%%%%%%%%%%%%%%%%%%%%%%%%%%%%%

%%%%%%%%%%%%%%%%%%%%The morphism with unstackable image words%%%%%%%%%%%%%
\section{The Morphism With Unstackable Image Words}

In a vain attempt to classify all of the overlap-free morphisms using the latin square morphism
\cite{tompkins}, we stumbled across the Leech square-free morphism in \cite{albook}. The following is the
Leech square-free morphism

\begin{singlespace}
\begin{eqnarray*}
&h&\\
0 & \mapsto & 0121021201210\\
1 & \mapsto & 1202102012021\\
2 & \mapsto & 2010210120102,
\end{eqnarray*}
\end{singlespace}

\noindent
which originally appeared in \cite{leech}. Noticing that this morphism was overlap-free put a hole in our
attempt to classify all of the overlap-free morphisms using Latin square morphisms. But on the other
hand, we now could potentially find another class of overlap-free morphisms that could be explained in a
better manner than with test-sets as in \cite{ric_testsets}.

Using the test-set result in Chapter 3, we found the following overlap-free morphisms on four letters

\begin{singlespace}
\begin{eqnarray*}
&f&\\
0 & \mapsto & 0123 1230 1 0321 3210\\
1 & \mapsto & 1230 2301 2 1032 0321\\
2 & \mapsto & 2301 3012 3 2103 1032\\
3 & \mapsto & 3012 0123 0 3210 2103,
\end{eqnarray*}
\end{singlespace}

\noindent
and

\begin{singlespace}
\begin{eqnarray*}
&g&\\
0 & \mapsto & 0123 0122121120 3210\\
1 & \mapsto & 1230 1300303301 0321\\
2 & \mapsto & 2301 2012331022 1032\\
3 & \mapsto & 3012 3011010013 2103.
\end{eqnarray*}
\end{singlespace}

\noindent
The morphism $g$ raised a considerable number of questions as to why it was overlap-free. It seemed to
avoid a considerable number of the techniques used in the proof For the Latin square morphisms. So the
natural question was: what does the morphism $g$ have in common with the Leech square-free
morphism that causes its overlap-freeness.

%%%%%%%%%%%%%%%%%%%%%%%%%%%%%%%%%%%%%%%%%%%%%%%%%%%%%%%%%%%%%%%%%%%%%%%%%%

%%%%%%%%%%%%%%%%%%%%%Definitions and Theorems%%%%%%%%%%%%%%%%%%%%%%%%%%%%%
\section{Definitions and Theorems}

The overlap-free morphisms displayed above are tied together with the following definition.

\begin{defn}
Let $h:\Sigma^*\to\Delta^*$ be an $n$-uniform morphism. We say that $h$ is a morphism with
unstackable image words if it
satisfies the
following properties:
\begin{itemize}
\item[(i)] $h(W)$ is overlap-free for all overlap-free words $W\in\Sigma^*$ with $|W|=3$.
\item[(ii)] For $a,b\in\Sigma$, and for all $V\in\Sigma^*$ such that $|V|\le\lfloor n/2\rfloor$,
\[h(a) = SV\quad\textrm{and}\quad h(b) = VU\]
if and only if $S$ is not a suffix of any image word of $h$ and $U$ is not a prefix of any image word of $h
$.
\end{itemize}
\label{def_5_1}
\end{defn}

We now prove a lemma that captures the combinatorial properties in the first
portion of Definition \ref{def_5_1}.

\begin{lem}
Let $\Sigma$ be an alphabet with more than one letter. Let $h:\Sigma^*\to\Delta^*$ be a morphism such
that $h(W)$ is overlap-free for all overlap-free $W\in
\Sigma^*$ with $|W| = 3$. We then have the following properties:
\begin{itemize}
\item[(i)] $h(a)$ is overlap-free for all $a\in\Sigma$.
\item[(ii)] $h(a)h(b)$ is overlap-free for all $a,b\in\Sigma$.
\item[(iii)] $h(a)$ and $h(b)$ do not begin or end with the same letter, whenever $a,b\in\Sigma$ and $a
\neq b$.
\end{itemize}
\label{lem_5_2}
\end{lem}

\begin{proof}
(i) Let us first state that the result does not apply when $|\Sigma|=1$ because there are no overlap-free
word of length three for this alphabet. For $|\Sigma|>1$ this result is clear because if we assume for a
contradiction that $h(a)$ contained an overlap for
any $a\in\Sigma$, then $h(bab)$, with $b\neq a$, would contain an overlap which contradicts our
assumption.

(ii) Similar to (i), if we assume that $h(a)h(b) =  h(ab)$ contained an overlap, then $h(aba)$ would contain
an overlap. Again this contradicts our assumption. We also must show that $h(aa)$ does not contain an
overlap. Assume for a contradiction that it does, and we quickly obtain our contradiction by observing that
then $h(aab)$ must contain an overlap.

(iii) Assume that for some $a,b\in\Sigma$, $h(a)$ and $h(b)$ begin with the same letter. Then,
$h(aab)$ would contain an overlap, which contradicts our assumption. The argument for $h(b)$ and
$h(b)$ ending with different letters is similar.
\end{proof}

\begin{thm}
Any morphism with unstackable image words is overlap-free.
\label{thm_5_3}
\end{thm}

\begin{proof}
We begin by assuming that $h$ is a morphism with unstackable image words with
$|h(a)|=n$ for all $a\in\Sigma$. We must show
that for all $W\in\Sigma^*$, $W$
is overlap-free if and only if $h(W)$ is overlap-free. We will begin with the easy direction first.

\subsection{The $\Leftarrow$ direction} Assume that $W=AcXcXcB$, so that we can argue by
contrapositive that $h(W)$ must also contain an
overlap. Notice that
\[h(W) = h(A)h(c)h(X)h(c)h(X)h(c)h(B).\]
Set $h(c) = dY$ where $d\in\Sigma$ and $Y\in\Sigma^*$, then $h(W) = h(A)dYh(X)dYh(x)dYh(B)$.
So then $h(W)$ contains the overlap $dYh(X)dYh(X)d$, and we are done with the first portion of our
argument.

\subsection{The $\Rightarrow$ direction}Conversely we will argue by contrapositive. We will assume that
$h(W)$ contains an overlap and show
that $W$ must also contain an overlap. So assume that for some $W\in\Sigma^*$ we have
\begin{displaymath}
h(W) = Ac_{j_0}Xc_{j_1}Xc_{j_2}B,
\end{displaymath}
where $c=c_{j_0}=c_{j_1}=c_{j_2}$. We use the 0, 1 and 2 to denote which $c$ we will refer to. Further,
the index $j_i$ will refer to which letter in the word $h(W)$ we are referring to, noting that we are indexing
beginning with 0.

We will proceed with two separate arguments. The first argument will be that it is not possible to write
$h(W)$ with $|cX|\not\equiv 0\pmod{n}$. The second argument will be that $W$ must contain an overlap
if $|cX|\equiv 0 \pmod{n}$.

\subsubsection{The $|cX|\not\equiv 0\pmod{n}$ case.} Notice that we must have the overlap in $h(W)$
contained
in $h(Z)$ where $|Z|>3$ is some subword of $W$. Otherwise we would be breaking hypothesis (i) in the
definition of Pooh morphisms.

We begin by setting
\[r_i\equiv j_i\pmod{n},\]
where $i\in\{0,1,2\}$ and $r_i\in\{0,1,\ldots,n-1\}$. We will argue first based on the number of tiles that the
overlap occurs in, and then by cases. When the overlap occurs over four tiles (noting that occurring over
three tiles contradicts the hypothesis), we will observe four cases. The cases are

\begin{singlespace}
\begin{align*}
r_0\le r_2<r_1,\\
r_2<r_0<r_1,\\
r_1<r_0 \le r_2,\\
r_1<r_2<r_0.\\
\end{align*}
\end{singlespace}

\noindent
We note that the cases $r_0<r_1<r_2$ and $r_2<r_1<r_0$ force the overlap to occur in a number other
than four tiles. When the overlap occurs in more than four tiles we will more simply consider the two
cases $r_0<r_1$ and $r_1<r_0$. Finally, we note the following relationship between $r_0$, $r_1$, and
$r_2$.
\begin{equation}
r_2\equiv 2r_1 - r_0 \pmod{n}.
\label{eq_5_1}
\end{equation}

Consider the notion of the tiling of a line segment. We will use this notion of tiling in application to working
with $h(W)$. The tiles we speak of are the image words of $h$. Note that all the image words must be of
the
same length $n$, this is crucial to our argument. For ease we will use $T_{s_i}$ with $i \in\{0,1,2\}$ to
denote the tile containing $c_{j_i}$. Note that $s_i$ is the number of the tile if we numbered them
starting with the first tile as $T_{0}$.

\emph{The overlap is contained in 4 tiles.} Let us consider the case where there is some subword of $W$,
say $Z$, with $|Z|=4$ and the overlap in $h(W)$ is contained in $h(Z)$. As in the argument for a Latin
square morphism to be overlap-free we will consider the word $h(Z)$ to be a line. We will draw in small
vertical lines to signify the edges of the tiles, and we will draw taller labeled vertical lines to signify the $c
$'s in the overlap.

\begin{figure}[h]
\centerline{
\begin{tikzpicture}
%
%Thesis: short r_2<r_0<r_1
%
%Main sequence lines
\draw (-3.1,1.0)--(6.9,1.0);
\draw (-6.8,-1.0)--(3.2,-1.0);
%Tile edges bottom line
\draw (-6.8,-1.2)--(-6.8,-.8);
\draw (-4.3,-1.2)--(-4.3,-.8);
\draw (-1.8,-1.2)--(-1.8,-.8);
\draw (.7,-1.2)--(.7,-.8);
\draw (3.2,-1.2)--(3.2,-.8);
%Tile edges upper line
\draw (-3.1,.8)--(-3.1,1.2);
\draw (-.6,.8)--(-.6,1.2);
\draw (1.9,.8)--(1.9,1.2);
\draw (4.4,.8)--(4.4,1.2);
\draw (6.9,.8)--(6.9,1.2);
%c's on upper line
\draw (-2.8,.6)--(-2.8,1.4);
\draw (.9,.6)--(.9,1.4);
\draw (4.6,.6)--(4.6,1.4);
%c's on bottom line
\draw (-2.8,-.6)--(-2.8,-1.4);
\draw (.9,-.6)--(.9,-1.4);
\draw (-6.5,-.6)--(-6.5,-1.4);
%write c's on upper line
\coordinate[label=above:$c_{j_0}$] (A) at (-2.8,1.4);
\coordinate[label=above:$c_{j_1}$] (B) at (.9,1.4);
\coordinate[label=above:$c_{j_2}$] (C) at (4.6,1.4);
%write c's on lower line
\coordinate[label=below:$c_{j_0}$] (a) at  (-6.5,-1.4);
\coordinate[label=below:$c_{j_1}$] (b) at (-2.8,-1.4);
\coordinate[label=below:$c_{j_2}$] (c) at  (.9,-1.4);
%Draw in lines for V
\draw[fill=white] (-1.8,1.07) rectangle (-.6,.93);
\draw[fill=white] (-1.8,-.93) rectangle (-.6,-1.07);
%Draw arrow for V's
\draw[<->](-1.2,-.6)--(-1.2,.6) ;
%Write in V
\coordinate[label=above:$V$] (Vu) at (-1.2,1.4);
\coordinate[label=below:$V$] (vd) at (-1.2,-1.4);
%Draw in U lines
\draw[fill=black] (-.6,1.07) rectangle (.7,.93);
\draw[fill=black] (-.6,-.93) rectangle (.7,-1.07);
%Draw U's
\coordinate[label=above:$U$] (Uu) at (.05,1.4);
\coordinate[label=below:$U$] (Ud) at (.05,-1.4);
%Write not arrow
\draw[<->](.05,.6)--(.05,-.6);
\end{tikzpicture}
}
\caption{The short overlap with $r_2<r_0<r_1$}\label{fig_5_1}
\end{figure}

In Figure \ref{fig_5_1}, we have taken $c_{j_0}Xc_{j_1}Xc_{j_2}$ and written it twice aligning $c_{j_0}Xc_{j_1}$
in the upper line with $c_{j_1}Xc_{j_2}$ in the lower line for the purpose of equating the terms through
the overlap. Figure \ref{fig_5_1} displays the case when $r_2<r_0<r_1$. We remark here that the case when
$r_2=r_0<r_1$ proceeds in the same manner.

Let $V$ to be the final $r_1-r_0$ letters in the tile $T_{s_0}$, as we have drawn in Figure \ref{fig_5_1}.  Similarly
we choose $U$ to be the first $r_1-r_2$ letters in $T_{s_1}$. Now equation (\ref{eq_5_1}) gives that in
the $r_2<r_0<r_1$ situation we have that $n-(r_1-r_0) = r_1-r_2$. Clearly then we must have $|U|=r_1-
r_2\le\lfloor n/2 \rfloor$ or $|V|=r_1-r_0\le\lfloor n/2 \rfloor$. In the case when $|V|\le\lfloor n/2\rfloor$ we
cannot equate $U$ with any prefix of a tile which leads to a contradiction. In the other case when $|U|\le
\lfloor n/2\rfloor$ we cannot equate $V$ with any suffix of a tile which leads to a contradiction. So this
case is not possible.

We now consider the case where $r_0<r_2<r_1$ as shown in Figure \ref{fig_5_a}.

\begin{figure}[h]
\centerline{
\begin{tikzpicture}
%
%Thesis: short r_0<r_2<r_1
%
%Main sequence lines
\draw (-4.5,1.0)--(5.5,1.0);
\draw (-8.7,-1.0)--(1.3,-1.0);
%Tile edges upper line
\draw (-4.5,.8)--(-4.5,1.2);
\draw (-2.0,.8)--(-2.0,1.2);
\draw (.5,.8)--(.5,1.2);
\draw (3.0,.8)--(3.0,1.2);
\draw (5.5,.8)--(5.5,1.2);
%Tile edges lower line
\draw (-8.7,-.8)--(-8.7,-1.2);
\draw (-6.2,-.8)--(-6.2,-1.2);
\draw (-3.7,-.8)--(-3.7,-1.2);
\draw (-1.2,-.8)--(-1.2,-1.2);
\draw (1.3,-.8)--(1.3,-1.2);
%write c's on upper line
\coordinate[label=above:$c_{j_0}$] (A) at (-4.2,1.4);
\coordinate[label=above:$c_{j_1}$] (B) at (0,1.4);
\coordinate[label=above:$c_{j_2}$] (C) at (4.2,1.4);
%write c's on lower line
\coordinate[label=below:$c_{j_0}$] (a) at  (-8.4,-1.4);
\coordinate[label=below:$c_{j_1}$] (b) at (-4.2,-1.4);
\coordinate[label=below:$c_{j_2}$] (c) at  (0,-1.4);
%Draw in c lines on upper line
\draw (-4.2,.6)--(-4.2,1.4);
\draw (0,.6)--(0,1.4);
\draw (4.2,.6)--(4.2,1.4);
%Draw in c lines on lower line
\draw (-8.4,-1.4)--(-8.4,-.6);
\draw (-4.2,-1.4)--(-4.2,-.6);
\draw (0,-1.4)--(0,-.6);
%Draw in lines for V
\draw[fill=black] (-2,1.07) rectangle (-1.2,.93);
\draw[fill=black] (-2,-.93) rectangle (-1.2,-1.07);
%Write in V
\coordinate[label=above:$U$] (Vu) at (-1.6,1.4);
\coordinate[label=below:$U$] (vd) at (-1.6,-1.4);
%Draw in U lines
\draw[fill=white] (-3.7,1.07) rectangle (-2.0,.93);
\draw[fill=white] (-3.7,-.93) rectangle (-2.0,-1.07);
%Draw U's
\coordinate[label=above:$V$] (Uu) at (-2.8,1.4);
\coordinate[label=below:$V$] (Ud) at (-2.8,-1.4);
%Draw in equating arrow
\draw[<->](-1.6,-.6)--(-1.6,.6);
%Draw in contradiction arrow
\draw[<->](-2.85,-.6)--(-2.85,.6);
\end{tikzpicture}
}
\caption{The short overlap with $r_0<r_2<r_1$}\label{fig_5_a}
\end{figure}

\noindent
In this case we choose $V$ to be the final $r_1-r_0$ letters in $T_{s_0}$, and we also pick $U$ to be the
first $r_1-r_2$ letters in $T_{s_1}$ as drawn in Figure \ref{fig_5_a}. Again we notice that $n-(r_1-r_0) = r_1-r_2$
so either $V\le\lfloor n/2\rfloor$ or $|U|\le\lfloor n/2\rfloor$, either of which is impossible. So we cannot
have this case occurring either.

\begin{figure}[h]
\centerline{
\begin{tikzpicture}
%
%Thesis: short r_1<r_2<r_0
%
%Draw the sequence lines
\draw (-6.5,-1.0)--(3.5,-1.0);
\draw (-3.0,1.0)--(7.0,1.0);
%Tile edges upper line
\draw (-3.0,.8)--(-3.0,1.2);
\draw (-.5,.8)--(-.5,1.2);
\draw (2.0,.8)--(2.0,1.2);
\draw (4.5,.8)--(4.5,1.2);
\draw (7.0,.8)--(7.0,1.2);
%Tile edges bottom line
\draw (-6.5,-1.2)--(-6.5,-.8);
\draw (-4.0,-1.2)--(-4.0,-.8);
\draw (-1.5,-1.2)--(-1.5,-.8);
\draw (1.0,-1.2)--(1.0,-.8);
\draw (3.5,-1.2)--(3.5,-.8);
%Lines for c's upper
\draw (-1.0,.6)--(-1.0,1.4);
\draw (2.5,.6)--(2.5,1.4);
\draw (6.0,.6)--(6.0,1.4);
%Lines for c's lower
\draw (-4.5,-1.4)--(-4.5,-.6);
\draw (-1.0,-1.4)--(-1.0,-.6);
\draw (2.5,-1.4)--(2.5,-.6);
%write c's on upper line
\coordinate[label=above:$c_{j_0}$] (A) at (-1.0,1.4);
\coordinate[label=above:$c_{j_1}$] (B) at (2.5,1.4);
\coordinate[label=above:$c_{j_2}$] (C) at (6.0,1.4);
%write c's on lower line
\coordinate[label=below:$c_{j_0}$] (a) at  (-4.5,-1.4);
\coordinate[label=below:$c_{j_1}$] (b) at (-1.0,-1.4);
\coordinate[label=below:$c_{j_2}$] (c) at  (2.5,-1.4);
%Draw in lines for V
\draw[fill=black] (1.0,1.07) rectangle (2,.93);
\draw[fill=black] (1.0,-.93) rectangle (2,-1.07);
%Write in V
\coordinate[label=above:$U$] (Vu) at (1.5,1.4);
\coordinate[label=below:$U$] (vd) at (1.5,-1.4);
%Draw in U lines
\draw[fill=white] (-.5,1.07) rectangle (1,.93);
\draw[fill=white] (-.5,-.93) rectangle (1,-1.07);
%Draw U's
\coordinate[label=above:$V$] (Uu) at (.25,1.4);
\coordinate[label=below:$V$] (Ud) at (.25,-1.4);
%Draw V arrow
\draw[<->](1.5,-.6)--(1.5,.6);
%Draw U arrow
\draw[<->](.25,-.6)--(.25,.6);
\end{tikzpicture}
}
\caption{The short overlap with $r_1<r_2<r_0$}\label{fig_5_2}
\end{figure}

We now consider the cases with $r_1<r_2\le r_0$ and $r_1<r_0<r_2$. Figure \ref{fig_5_2} gives the situation
when $r_1<r_2<r_0$ (note that the case when $r_1<r_2=r_0$ is similar, and the same applies to the
arguments above). Notice that in both of these cases we have that $n-(r_2-r_1) = r_0-r_1$.



For the case depicted in Figure \ref{fig_5_2}, $r_1<r_2<r_0$, we assume that $V$ is the final $r_0-r_1$ letters in
$T_{s_1}$, and we also assume that $U$ is the first $r_2-r_1$ letters in $T_{s_2}$. Now either $|U|\le
\lfloor n/
2\rfloor$ or $|V|\le\lfloor n/2\rfloor$. In either case we have a contradiction.

Now we consider the case where $r_1<r_0<r_2$, which is displayed in Figure \ref{fig_5_b}.
\begin{figure}[h]
\centerline{
\begin{tikzpicture}
%
%Thesis: short r_1<r_0<r_2
%
%Draw upper and lower lines
\draw (-5.3,1.0)--(4.8,1.0);
\draw (-9.2,-1.0)--(.8,-1.0);
%Draw tile edges on upper line
\draw (-5.3,.8)--(-5.3,1.2);
\draw (-2.8,.8)--(-2.8,1.2);
\draw (-.3,.8)--(-.3,1.2);
\draw (2.2,.8)--(2.2,1.2);
\draw (4.7,.8)--(4.7,1.2);
%Draw tile edges on lower line
\draw (-9.2,-1.2)--(-9.2,-.8);
\draw (-6.7,-1.2)--(-6.7,-.8);
\draw (-4.2,-1.2)--(-4.2,-.8);
\draw (-1.7,-1.2)--(-1.7,-.8);
\draw (.8,-1.2)--(.8,-.8);
%Draw the c lines on upper line
\draw (-3.9,.6)--(-3.9,1.4);
\draw (0,.6)--(0,1.4);
\draw (3.9,.6)--(3.9,1.4);
%write c's on upper line
\coordinate[label=above:$c_{j_0}$] (A) at (-3.9,1.4);
\coordinate[label=above:$c_{j_1}$] (B) at (0,1.4);
\coordinate[label=above:$c_{j_2}$] (C) at (3.9,1.4);
%write c's on lower line
\coordinate[label=below:$c_{j_0}$] (a) at  (-7.8,-1.4);
\coordinate[label=below:$c_{j_1}$] (b) at (-3.9,-1.4);
\coordinate[label=below:$c_{j_2}$] (c) at  (0,-1.4);
%Draw the c lines on the lower line
\draw (-7.8,-1.4)--(-7.8,-.6);
\draw (-3.9,-1.4)--(-3.9,-.6);
\draw (0,-1.4)--(0,-.6);
%Draw in lines for V
\draw[fill=white] (-2.8,1.07) rectangle (-1.7,.93);
\draw[fill=white] (-1.7,-.93) rectangle (-2.8,-1.07);
%Write in V
\coordinate[label=above:$V$] (Vu) at (-2.25,1.4);
\coordinate[label=below:$V$] (vd) at (-2.25,-1.4);
%Draw in U lines
\draw[fill=black] (-1.7,1.07) rectangle (-.3,.93);
\draw[fill=black] (-.3,-.93) rectangle (-1.7,-1.07);
%Draw U's
\coordinate[label=above:$U$] (Uu) at (-.7,1.4);
\coordinate[label=below:$U$] (Ud) at (-.7,-1.4);
%Draw equating arrow
\draw[<->](-2.25,.6)--(-2.25,-.6);
%Draw contradiction arrow
\draw[<->](-1.0,.6)--(-1.0,-.6);
\end{tikzpicture}
}
\caption{The short overlap with $r_1<r_0<r_2$}\label{fig_5_b}
\end{figure}

In the case displayed here in Figure \ref{fig_5_b}, we again assume that $V$ occurs in the final $r_0-r_1$ letters
of $T_{s_1}$, and we also assume that $U$ occurs in the first $r_2-r_1$ letters of $T_{s_2}$. We then
have that either $|V|\le\lfloor n/2\rfloor$ or that $|U|\le\lfloor n/2\rfloor$. Either case is a contradiction. So
we cannot have our overlap occurring in four tiles. Thus, we consider the case when the overlap occurs
in more than four tiles.

We also not that if we are in the case when the overlap occurs in five tiles, the same arguments hold.

\emph{The overlap is contained in more than four tiles.} We will look at the cases with $r_0<r_1$ and
$r_1<r_0$, and we will only look at the beginning of the overlap. So we consider Figure \ref{fig_5_3} for the case
when $r_0<r_1$.

\begin{figure}[h]
\centerline{
\begin{tikzpicture}
%Start by drawing the sequence lines
\draw (-6.0,-1.0)--(3.5,-1.0);
\draw (-4.0,1.0)--(2.5,1.0);
%Put the dots on the end
\coordinate[label=right:$\bullet$] (e) at (3,0);
\coordinate[label=right:$\bullet$] (f) at (3.5,0);
\coordinate[label=right:$\bullet$] (g) at (4,0);
%Draw in the tile edges on the upper line
\draw (-4.0,.8)--(-4.0,1.2);
\draw (-1.0,.8)--(-1.0,1.2);
\draw (2.0,.8)--(2.0,1.2);
%Draw the tile edges on the lower line
\draw (-6.0,-.8)--(-6.0,-1.2);
\draw (-3.0,-.8)--(-3.0,-1.2);
\draw (0,-.8)--(0,-1.2);
\draw (3.0,-.8)--(3.0,-1.2);
%Draw in the line and the t_{j_0}
\coordinate[label=above:$c_{j_0}$] (A) at (-3.5,1.4);
\draw (-3.5,.6)--(-3.5,1.4);
%Draw in the line and the t_{j_1}
\coordinate[label=below:$c_{j_1}$] (a) at (-3.5,-1.4);
\draw (-3.5,-.6)--(-3.5,-1.4);
%Draw in V line
\draw[fill=white] (-3,1.07) rectangle (-1,.93);
\draw[fill=white] (-1,-.93) rectangle (-3,-1.07);
%Draw V
\coordinate[label=above:$V$] (Vu) at (-2,1.4);
\coordinate[label=below:$V$] (Vd) at (-2,-1.4);
%Draw arrow
\draw[<->](-2.0,-.6)--(-2.0,.6);
%Draw U line
\draw[fill=black] (-1,1.07) rectangle (0,.93);
\draw[fill=black] (0,-.93) rectangle (-1,-1.07);
%Draw U
\coordinate[label=above:$U$] (Uu) at (-.5, 1.4);
\coordinate[label=below:$U$] (Ud) at (-.5,-1.4);
%Draw arrow
\draw[<->](-.5,.6)--(-.5,-.6);
%Contradiction arrow
\draw[<-](1.0,-.6)--(1.0,.6);
\draw[-](.9,-.1)--(1.1,.1);
\end{tikzpicture}
}
\caption{The long overlap with $r_0<r_1$ and $r_1-r_0\ge\lfloor n/2\rfloor$}\label{fig_5_3}
\end{figure}
We will consider the case with $r_1-r_0\ge\lfloor n/2\rfloor$, as we will cover the logic behind the
argument for $r_1-r_0\le\lfloor n/2\rfloor$ in Figure \ref{fig_5_4}.

Let $V$ be the final $r_1-r_0$ letters in $T_{s_0}$, then equating yields $V$ as the beginning $r_0-r_1$
letters of $T_{s_1+1}$. Since $|V|\ge \lfloor n/2\rfloor$ we can equate the suffix of $T_{s_1+1}$, call it $U$
(which is labeled with a dotted line in Figure \ref{fig_5_3}),
with $T_{s_0+1}$. So we have $T_{s_1+1} = VU$. Similarly we can set $S\in\Delta^*$ such that
$T_{s_0+1} = US$. Notice now that $|U| = n-(r_1-r_0)\le\lfloor n/2\rfloor$. Thus $S$ cannot begin any
image word of $h$ so the overlap is impossible.

A note for the case when $r_1-r_0\le\lfloor n/2\rfloor$. In this case $|V|\le\lfloor n/2\rfloor$ and we would
not be able to equate $U$.

\begin{figure}[h]
\centerline{
\begin{tikzpicture}
%Draw the upper line
\draw (-6.0,1.0)--(.5,1.0) ;
%Draw lower line
\draw (-4.5,-1.0)--(2.0,-1.0) ;
%Draw dots at end of lines
\coordinate[label=right:$\bullet$] (e) at (1.5,0);
\coordinate[label=right:$\bullet$] (f) at (2,0);
\coordinate[label=right:$\bullet$] (g) at (2.5,0);
%Draw the upper tile edges
\draw (-6.0,.8)--(-6.0,1.2) ;
\draw (-3.0,.8)--(-3.0,1.2) ;
\draw (0,.8)--(0,1.2);
%Draw the lower tile edges
\draw (-4.5,-1.2)--(-4.5,-.8) ;
\draw (-1.5,-1.2)--(-1.5,-.8) ;
\draw (1.5,-1.2)--(1.5,-.8);
%Draw the line for c_{j_0}
\draw (-4.0,.6)--(-4.0,1.4) ;
\coordinate[label=above:$c_{j_0}$] (A) at (-4.0,1.4);
%Draw the line for c_{j_1}
\draw (-4.0,-1.4)--(-4.0,-.6) ;
\coordinate[label=below:$c_{j_1}$] (b) at (-4.0,-1.4);
%Draw \dir2{-} lines for V
\draw[fill=white] (-3,1.07) rectangle (-1.5,.93);
\draw[fill=white] (-3,-.93) rectangle (-1.5,-1.07);
%Draw V
\coordinate[label=above:$V$] (Vu) at (-2.25,1.4);
\coordinate[label=below:$V$] (Vd) at (-2.25,-1.4);
%Draw the equating arrow
\draw[<->](-2.25,-.6)--(-2.25,.6);
%Draw contradiction arrow
\draw[<->](-.75,-.6)--(-.75,.6);
\draw (-.85,-.1)--(-.65,.1);
\end{tikzpicture}
}
\caption{The long overlap with $r_1<r_0$ and $r_0-r_1\le\lfloor n/2\rfloor$}\label{fig_5_4}
\end{figure}
Figure \ref{fig_5_4} gives the case when $r_1<r_0$ with $r_0-r_1\le\lfloor n/2\rfloor$. In a similar manner to the
case where $r_0<r_1$ we pick $V$ to be the suffix of $r_0-r_1$ letters in $T_{s_1}$. Now we can
possibly find an image word $T_{s_0+1} = VU$ for some $U\in\Delta^*$. But because $|V| = r_0-r_1\le
\lfloor n/2\rfloor$, $U$ cannot be the prefix of any image word, so this formulation of the overlap is
impossible.

So we see that in order for $h(W)$ to contain an overlap, it must be one so that $|cX|\equiv 0 \pmod{n}$.

\subsubsection{The $|cX|\equiv 0 \pmod{n}$ case.} From Lemma \ref{lem_5_2} we know that the
beginning letters and ending letter for each image word in $h$ must be distinct. Further we know that the
suffix of $T_{s_1}$ must be identical to the suffix of $T_{s_0}$ as $r_0=r_1$. This implies that
$T_{s_0}=T_{s_1}$. Similarly $T_{s_1} = T_{s_2}$.

Pick $d$ to be the letter such that $h(d) = T_{s_0} = T_{s_1} = T_{s_2} = T$. Further because
\[h(W) = Ac_{j_0}Xc_{j_1}Xc_{j_2}XB,\]
we can find subwords $C,D,Y$ of $W$ so that
\[h(W) = h(CdYdYdD) = AcXcXcB.\]
Now we must have that $W = CdYdYdD$ which contains an overlap. Thus we are done.


\end{proof}

%%%%%%%%%%%%%%%%%%%%%%%%%%%%%%%%%%%%%%%%%%%%%%%%%%%%%%%%%%%%%%%%%%%%%%%%%%

%%%%%%%%%%%%%%%%%%%%The square-free unstackable morphism%%%%%%%%%%%%%%%%%%
\section{The Square-Free Adaptation}

Similarly to the definition of the overlap-free morphisms with unstackable image words we can define
square-free morphisms with unstackable image words in the
following manner.

\begin{defn}
Let $h:\Sigma^*\to\Delta^*$ be an $n$-uniform morphism. We call $h$ a square-free morphism with
unstackable image words if it
satisfies the following properties:
\begin{itemize}
\item[(i)] $h(W)$ is square-free for all square-free words $W\in\Sigma^*$ with $|W| = 3$
\item[(ii)] $h(a)$ and $h(b)$ do not begin or end with the same letter for all $a,b\in\Sigma$ with $a\neq b$.
\item[(iii)] For $a,b\in\Sigma$, and for all $V\in\Sigma^*$ such that $|V|\le\lfloor n/2\rfloor$,
\[h(a) = SV\quad\textrm{and}\quad h(b) = VU\]
if and only if $S$ is not a suffix of any image word of $h$ and $U$ is not a prefix of any image word of $h
$.
\end{itemize}
\label{defn_5_4}
\end{defn}

Because we cannot consider words like $aab$ to put into $h$, we must add property (ii) in Definition \ref{defn_5_4}
so that we can use a similar preimage argument in the final portion of the argument. Thus we have the following
theorem.

\begin{thm}
Any square-free morphism with unstackable image words is square-free.
\label{thm_5_5}
\end{thm}

\begin{proof}
Assume that $h$ is a square-free morphism with unstackable image words such that $|h(a)| = n$ for all $a\in\Sigma$. We must
show that for some $W\in\Sigma^*$, $W$ is square-free if and only if $h(W)$ is square-free. We will begin
with the easy direction.

\subsection{The $\Leftarrow$ direction} We will proceed by contrapositive. So assume that $W = AXXB$,
where $X\in\Sigma^+$ and $A,B\in\Sigma^*$. Write
\[h(W) = h(AXXB) = h(A)h(X)h(X)h(B),\]
which contains the square $h(X)h(X)$. So we are done with this direction.

\subsection{The $\Rightarrow$ direction} Again we proceed by arguing the contrapositive. So we assume
that
\begin{equation}
h(W) = Ac_{j_0}Xd_{i_0}c_{j_1}Xd_{i_1}B,
\label{eq_5_3}
\end{equation}
where $c=c_{j_0}=c_{j_1}\in\Sigma$, $d = d_{i_0} = d_{i_1}\in\Sigma$ and $A,X,B\in\Sigma^*$. Note that
we are using $c_{j_0}$ and $c_{j_1}$ so that we can mark the beginning of the square, and similarly for
the $d$'s and the end of the square.

There are two cases to consider here $|cXd|\not\equiv 0 \pmod{n}$ and $|cXd|\equiv 0 \pmod{n}$. We
show
that it is impossible for $|cXd|\not\equiv 0 \pmod{n}$ in an analogous manner as in Theorem
\ref{thm_5_3},
as seen in section 5.1.2.1.

So assume that $|cXd|\equiv 0 \pmod{n}$. From the definition of the pooh square-free morphism we know
that each image word for $h$ must begin and end with distinct letters. Further we know the suffix of the
tile containing $c_{j_0}$ must be identical to the suffix of the tile containing $c_{j_1}$. Thus they are the
same image word, call it $h(z)$ for some $z\in\Sigma$. So because
\[h(W) = AcXcXB,\]
we can find subwords $C,D,Y$ of $W$ such that
\[h(W) = h(CzYzYD) = AcXcXB.\]
Thus we have that $W = AzXzXB$ which contains a square, and we are done.
\end{proof}
%%%%%%%%%%%%%%%%%%%%%%%%%%%%%%%%%%%%%%%%%%%%%%%%%%%%%%%%%%%%%%%%%%%%%%%%%%

%%%%%%%%%%%%%%%%%%%%%Acknowlegements%%%%%%%%%%%%%%%%%%%%%%%%%%%%%%%%%%%%%%
\section{Acknowlegments}

%%%%%%%%%%%%%%%%%%%%%%%%%%%%%%%%%%%%%%%%%%%%%%%%%%%%%%%%%%%%%%%%%%%%%%%%%%


\begin{thebibliography}{99}
 \bibitem{al} J.-P. Allouche and J. Shallit. \emph{Automatic Sequences:
 Theory, Applications, Generalizations}. Cambridge Press, Cambridge,
 UK. (2003) 1-17.
 \bibitem{al2} J.-P. Allouche and J. Shallit. The ubiquitious
 Prouhet-Thue-Morse sequence. In C. Ding, T. Helleseth and H.
 Niederreiter, eds., \emph{Sequences and Their Applications,
 Proceedings of SETA'98}, Springer-Verlag, 1999, 1-16.
 \bibitem{al3} J.-P. Allouche and J. Shallit. Sums of Digits,
 Overlaps, and Palindromes. \emph{Discrete Mathematics and
 Theoretical Computer Science} \textbf{4}, 2000, 001-010.
 \bibitem{seebold} J. Berstel and P. S\'e\'ebold. A
 characterization of overlap-free morpisms. \emph{Discrete Appl. Math.}
 \textbf{46}, no. 3, 1993, 275-281.
 \bibitem{de} J. D\'enes and A.D. Keedwell. \emph{Latin Squares and
 Their Applications}. Academic Press Inc., New York, New York and
 London, UK. (1974), 1-27.
 \bibitem{fr} A. Frid. Overlap-Free Symmetric D0L words.
 \emph{Discrete Mathematics and Theoretical Computer Science}
 \textbf{4}, 2001, 357-362.
 \bibitem{lothaire} Lothaire, M. \emph{Algebraic combinatorics on
 words}, Encyclopedia of Mathematics and its Applications, vol. 90,
 Cambridge University Press, Cambridge. 2002.
 \bibitem{maor} Maor, E., \emph{To Infinity and Beyond,} Princeton: Princeton Univeristy Press, 1987
 \bibitem{richomme} G. Richomme and F. Wlazinski. Overlap-free
 morphisms and finite test-sets. \emph{Discrete Appl. Math.} \textbf{143}, no. 1-3, 2004
 , 92-109.
 \bibitem{th} A. Thue, \"Uber die gegenseitige Lage gleicher Teile
 gewisser Zeichenreihen, \emph{Norske vid. Selsk. Skr. Mat. Nat. Kl.}
 \textbf{1} (1912), 1-67. Reprinted in "Selected mathematical papers of Axel Thue," T.
 Nagell, ed., Universitetsforlaget, Oslo, 1977, 413-478.
 \bibitem{tompkins} Tompkins, C. Robinson. Latin square Thue-Morse sequences
 are overlap-free.\emph{Discrete Mathematics and Theoretical Computer
 Science}. \textbf{9}, 2007, 239-246.
\end{thebibliography}


\end{onehalfspace}
\end{document}
